% !TeX root = main.tex
\section{Experimental Results}

This section presents the measured performance of the gPCD framework
for randomly distributed particle configurations of increasing size.
Unless otherwise noted, benchmark cases maintain an approximately
constant collision fraction while varying the total particle count.

To evaluate memory-locality effects, two occupancy configurations were
considered. In the first configuration, cell occupancy data remained
sequentially ordered in memory. In the second configuration, occupancy
data were randomly distributed throughout the cell array.

\Figo{A_PQBS_V_PQBR} compares execution time as a function of
particle count for sequential and randomly ordered occupancy data.
Sequential ordering consistently produced lower execution times
throughout the evaluated range. The figure also presents the
execution-time difference, defined as the random-ordering execution
time minus the sequential-ordering execution time.

\begin{table}[t]
\caption{Performance summary for randomly distributed particle configurations.}
\label{tab:performance_summary}
\centering
\begin{tabular}{lr}
\toprule
\textbf{Metric} & \textbf{Value} \
\midrule
Maximum Particle Count & \APQBRTTABLEALLMaxParticles \
Maximum Collision Count & \APQBRTTABLEALLMaxCollisions \
Peak Throughput (particles/s) & \APQBRTTABLEALLPeakThroughput \
Peak Frame Rate (FPS) & \APQBRTTABLEALLPeakFPS \
Largest-Case Compute Time (ms) & \APQBRTTABLEALLMaxComputems \
Largest-Case Graphics Time (ms) & \APQBRTTABLEALLMaxGraphicsms \
Largest-Case Total Time (ms) & \APQBRTTABLEALLMaxTotalms \
\bottomrule
\end{tabular}
\end{table}


Both configurations exhibit linear scaling behavior; however, the
execution-time difference increases with particle count. A linear
regression of the difference curve produced a slope of
\PPBQDiffSlope\ ms per million particles, indicating that the cost of
random memory access grows proportionally with problem size. Since the
two configurations differ only in the ordering of occupancy data, the
observed performance gap is attributed primarily to memory-locality
effects.

These results indicate that memory locality remains an important
contributor to overall framework performance despite the GPU-resident
nature of the implementation. Although memory locality does not alter
the asymptotic complexity of the method, the increasing performance gap
observed with particle count demonstrates its influence on practical
execution time.

\begin{table}[t]
\caption{Performance summary for randomly distributed particle configurations.}
\label{tab:performance_summary}
\centering
\begin{tabular}{lr}
\toprule
\textbf{Metric} & \textbf{Value}\\
\midrule
Maximum Particle Count & 10,474,496 \\ 
Maximum Collision Count & 5,237,248 \\ 
Peak Throughput (million particles/s) & 132 \\ 
Peak Frame Rate (FPS) & 13 \\ 
Largest-Case Compute Time (ms) & 65.72 \\ 
Largest-Case Graphics Time (ms) & 13.44 \\ 
Largest-Case Total Time (ms) & 79.16 \\ 
\bottomrule
\end{tabular}
\end{table}


\Tigo{performance_summary} summarizes the principal performance metrics
obtained during evaluation of the gPCD framework, including the largest
benchmark configuration and the maximum throughput achieved during
testing.

The framework successfully processed up to
\numprint{\APQBRTTABLEALLbka{}} particles and achieved a peak throughput
of \SUMMARYTHROUGHPUT{} million particles/s. Table~\ref{tab:performance_summary}
reports representative execution statistics for the largest benchmark
configuration.

Table~\ref{tab:A_PQBRT_TABLE_ALL} provides the benchmark measurements
used throughout the subsequent throughput, time-per-particle, and
scaling analyses.

Figure~\ref{fig:A_PQBR_THROUGHPUT_AND_TPP} presents throughput and
time-per-particle measurements across the evaluated particle range.
These metrics are used to examine resource utilization, large-scale
performance, and scaling behavior.


\begingroup
\begin{figure*}[ht]
\centering
\hspace{-.05in}
	\begin{subfigure}[b]{0.32\textwidth}
		\includegraphics[width=\textwidth]{A_PQBS_V_PQBR.png}
		\subcaption[]{random v sequential}
		\label{fig:A_PQBS_V_PQBR}
	\end{subfigure}
\hspace{-.05in}
	\begin{subfigure}[b]{0.32\textwidth}
		\includegraphics[width=\textwidth]{A_PQBR_THROUGH_PUT_ALL.png}
		\subcaption[]{throughput}
		\label{fig:A_PQBR_THROUGH_PUT_ALL}
	\end{subfigure}
\hspace{-.05in}
	\begin{subfigure}[b]{0.32\textwidth}
		\includegraphics[width=\textwidth]{A_PQBR_TIME_PERP_ALL.png}
		\subcaption[]{time per particle}
		\label{fig:A_PQBR_TIME_PERP_ALL}
	\end{subfigure}
\hspace{-.05in}
\caption[TITLE:PQBR Performance Plots]{(a) Comparison of randomly ordered particle occupancy (red squares) and
sequentially ordered particle occupancy (blue circles). Sequential
ordering consistently produces lower execution times throughout the
evaluated particle range. At the largest benchmark configuration,
random ordering increases execution time by approximately
\nc{Locality Penalty} relative to sequential ordering, highlighting
the importance of memory locality despite identical algorithmic
complexity.
(b) Log--log plot of total throughput (red), compute throughput (blue),
and graphics throughput (black) as a function of particle count. The
logarithmic axes permit visualization of performance trends across
nearly three orders of magnitude in particle count. Dashed vertical
lines indicate observed transition regions within the measured
throughput behavior.

A log-log plot of throughput for graphics, compute, and total. The saturated point plateau's after $N=$\numprint{\APQBRTIMEPERPALLstartval{}} particles shown in the plot as a black vertical dashed line. 

(c) time per particle for graphics, compute, total starting at, N=\numprint{\APQBRTIMEPERPALLstartval{}}, for the start of the heavily saturated range. Graphics stabilizes at $\APQBRTIMEPERPALLGraphicsstabilzes{}  \pm \APQBRTIMEPERPALLGraphicsstdval{} $ $\sfrac{ns}{particle}$, compute at $\APQBRTIMEPERPALLComputestabilzes{} \pm \APQBRTIMEPERPALLComputestdval{}$ $\sfrac{ns}{particle}$ for a total stabilization at $\APQBRTIMEPERPALLTotalstabilzes{} \pm \APQBRTIMEPERPALLTotalstdval{}$ $\sfrac{ns}{particle}$. }
\Description[TITLE:PQBR Performance Plots]{(a) A plot of randomly ordered particles performance (red squares) versus
sequentially ordered particles (blue circles). The random performance 
increases disproportionately with sequential performance as the number of 
particles increase.
(b) A log-log plot of throughput for graphics, compute, and total. The saturated point plateau's after $N=$\numprint{\APQBRTIMEPERPALLstartval{}} particles shown in the plot as a black vertical dashed line. (c) time per particle for graphics, compute, total starting at, N=\numprint{\APQBRTIMEPERPALLstartval{}}, for the start of the heavily saturated range. Graphics stabilizes at $\APQBRTIMEPERPALLGraphicsstabilzes{}  \pm \APQBRTIMEPERPALLGraphicsstdval{} $ $\sfrac{ns}{particle}$, compute at $\APQBRTIMEPERPALLComputestabilzes{} \pm \APQBRTIMEPERPALLComputestdval{}$ $\sfrac{ns}{particle}$ for a total stabilization at $\APQBRTIMEPERPALLTotalstabilzes{} \pm \APQBRTIMEPERPALLTotalstdval{}$ $\sfrac{ns}{particle}$. }
\label{fig:A_PQBR_THROUGHPUT_AND_TPP}
\end{figure*}
\endgroup

At low particle counts, throughput increases rapidly as available GPU
resources become more fully utilized. This \emph{utilization-limited region}
extends to approximately \numprint{100000} particles. Between
approximately \numprint{100000} and \numprint{2000000} particles,
throughput enters a \emph{throughput-transition regime} in which increasing occupancy and
collision workload begin to influence performance. Beyond
approximately \numprint{2000000} particles, the framework reaches a
\emph{stable large-scale regime}. Although throughput continues to
decrease gradually with increasing particle count, the corresponding
time-per-particle curves become nearly constant, indicating an
approximately constant average processing cost per particle.
Consequently, particle counts exceeding \numprint{2000000} are used for
the time-per-particle analysis reported in
Figure~\ref{fig:A_PQBR_THROUGHPUT_AND_TPP}(c).

\begingroup
\begin{figure*}[ht]
\centering
\hspace{-.1in}
	\begin{subfigure}[b]{0.32\textwidth}
		\includegraphics[width=\textwidth]{A_PQBR_LOGLOG_REGRESSgms.png}
		\subcaption[]{Graphics}
		\label{fig:A_PQBR_LOGLOG_REGRESSgms}
	\end{subfigure}
\hspace{-.1in}
	\begin{subfigure}[b]{0.32\textwidth}
		\includegraphics[width=\textwidth]{A_PQBR_LOGLOG_REGRESScms.png}
		\subcaption[]{Compute}
		\label{fig:A_PQBR_LOGLOG_REGRESScms}
	\end{subfigure}
\hspace{-.1in}
	\begin{subfigure}[b]{0.32\textwidth}
		\includegraphics[width=\textwidth]{A_PQBR_LOGLOG_REGRESStot.png}
		\subcaption[]{Total}
		\label{fig:A_PQBR_LOGLOG_REGRESStot}
	\end{subfigure}
\hspace{-.1in}
\caption[TITLE:PQBR LOG-LOG Regression]{Log--log regression of execution time $t_n$ versus particle count $N$. The measured data (circles) are fit by a straight line (solid), corresponding to a power-law model 
$t_n = \alpha N^{b}$. The fitted slope of $b \approx 1$ confirms near-linear 
scaling in practice, despite the $O(N^{2})$ theoretical bound of collision 
detection. The regression achieved $R^{2} > 0.99$ for compute (\APQBRLOGLOGREGRESScmsrtwosat{}) and graphics (\APQBRLOGLOGREGRESSgmsrtwosat{}) indicating that the throughput follows a consistent scaling law, with deviations only at small 
$N$ due to fixed overheads before GPU saturation is reached at 
$N \approx 10^{6}$.}
\Description[TITLE:PQBR LOG-LOG Regression]{Log--log regression of execution time $t_n$ versus particle count $N$. The measured data (circles) are fit by a straight line (solid), corresponding to a power-law model 
$t_n = \alpha N^{b}$. The fitted slope of $b \approx 1$ confirms near-linear 
scaling in practice, despite the $O(N^{2})$ theoretical bound of collision 
detection. The regression achieved $R^{2} > 0.99$ for compute (\APQBRLOGLOGREGRESScmsrtwosat{}) and graphics (\APQBRLOGLOGREGRESSgmsrtwosat{}) indicating that the throughput follows a consistent scaling law, with deviations only at small 
$N$ due to fixed overheads before GPU saturation is reached at 
$N \approx 10^{6}$.}
\label{fig:A_PQBR_LOGLOG_REGRESS}
\end{figure*}
\endgroup

Figure~\ref{fig:A_PQBR_LOGLOG_REGRESS} presents log--log regressions of
graphics, compute, and total execution times. A power-law model of the
form

\begin{equation}
t = aN^b
\end{equation}

was fitted to the measured data, where (t) is execution time, (N)
is particle count, (a) is a proportionality constant, and (b) is
the scaling exponent.

The measured scaling exponents were
$(b_g=\APQBRLOGLOGREGRESSgmsasat{})$,
$(b_c=\APQBRLOGLOGREGRESScmsasat{})$, and
$(b_t=\APQBRLOGLOGREGRESStotasat{})$.
All three exponents remain close to unity, indicating that execution
time increases approximately linearly with particle count throughout
the evaluated range.

\begingroup
\begin{figure*}[ht]
\centering
\hspace{-.1in}
	\begin{subfigure}[b]{0.32\textwidth}
		\includegraphics[width=\textwidth]{A_PBQR_CURVATUREgms.png}
		\subcaption[]{Compute}
		\label{fig:A_PBQR_CURVATUREgms}
	\end{subfigure}
\hspace{-.1in}
	\begin{subfigure}[b]{0.32\textwidth}
		\includegraphics[width=\textwidth]{A_PBQR_CURVATUREcms.png}
		\subcaption[]{Graphics}
		\label{fig:A_PBQR_CURVATUREcms}
	\end{subfigure}
\hspace{-.1in}
	\begin{subfigure}[b]{0.32\textwidth}
		\includegraphics[width=\textwidth]{A_PBQR_CURVATUREtot.png}
		\subcaption[]{Total}
		\label{fig:A_PBQR_CURVATUREtot}
	\end{subfigure}
\hspace{-.1in}
\caption[TITLE:PQB Curvature Bound Analysis]{Plots of quadratic fit lines for the saturated range of \numprint{\APBQRCURVATUREStartRowParticlesVal{}} to \numprint{\APBQRCURVATUREgmsMaxParticles{}} particles. The curvature for the graphics pipeline is \APBQRCURVATUREgmsa{}$N^2$. The curvature for the compute pipeline is \APBQRCURVATUREcmsa{}$N^2$ indicating a negative curvature. The total  curvature is \APBQRCURVATUREtota{}$N^2$. The total quadratic contribution is \APBQRCURVATUREtotQuadraticContributionpct{} percent.}
\Description[TITLE:PQB Curvature Bound Analysis]{Plots of quadratic fit lines for the saturated range of \numprint{\APBQRCURVATUREStartRowParticlesVal{}} to \numprint{\APBQRCURVATUREgmsMaxParticles{}} particles. The curvature for the graphics pipeline is \APBQRCURVATUREgmsa{}$N^2$. The curvature for the compute pipeline is \APBQRCURVATUREcmsa{}$N^2$ indicating a negative curvature. The total  curvature is \APBQRCURVATUREtota{}$N^2$. The total quadratic contribution is \APBQRCURVATUREtotQuadraticContributionpct{} percent.}
\label{fig:A_PBQR_CURVATURE}
\end{figure*}
\endgroup

While log--log regression quantifies the overall scaling exponent,
curvature analysis provides a complementary measure of deviation from
ideal linear behavior. Quadratic fits were therefore applied to the
graphics, compute, and total execution-time measurements.

Figure~\ref{fig:A_PBQR_CURVATURE} presents the resulting quadratic fits
and associated curvature bounds. In all cases, the fitted curves remain
close to the measured data, indicating only minor departures from
linear behavior. The graphics and total execution times exhibit
particularly small curvature, while the compute measurements display
slightly larger but still bounded deviations.

Combined with the log--log scaling exponents, these results provide
additional evidence that the gPCD framework maintains near-linear
scaling throughout the evaluated range.


